\section*{Notch Bandpass Sottrattore}

\noindent\colorbox{orange}{
    \begin{minipage}{1.0\linewidth}
      \paragraph{Richiesta (a)}
      Sulla base di semplici argomenti di natura qualitativa (ovvero senza effettuare il calcolo
      della funzione di trasferimento, come richiesto al punto successivo), mostrare che il
      blocco 1 del circuito svolga effettivamente la funzione di filtro passa-banda,
      discutendone la risposta ai limiti del dominio delle frequenze reali.
    \end{minipage}
}
\paragraph{Risposta (a)}
\begin{itemize}
\item
  Per $f \to 0$ i rami dei condensatori $C_1$ e $C_2$ si aprono, per impedenza d'ingresso
  ``elevata'' la resistenza $R_3$ \'e trascurabile per cui
  \begin{equation*}
    V_- = V_{BP}
  \end{equation*}
  e nella regione di operazione lineare dell'op-amp
  \begin{equation*}
    V_{BP} = A_d(0 - V_{BP}) \implies \colorbox{blu}{\boxed{V_{BP} = \frac{0}{1 + A_d} = 0}}
  \end{equation*}
\item
  Per $f \to \infty$ il ramo del condensatore $C_2$ si ``corto-circuita'', la tensione
  al nodo A, di conseguenza la tensione nell'ingresso invertente dell'op-amp
  \begin{equation*}
    V_- = V_{BP}
  \end{equation*}
  nella regione di operazione lineare dell'op-amp
  \begin{equation*}
    V_{BP} = A_d(0 - V_{BP}) \implies \colorbox{blu}{\boxed{V_{BP} = \frac{0}{1 + A_d} = 0}}
  \end{equation*}

\end{itemize}

\noindent\colorbox{orange}{
  \begin{minipage}{1.0\linewidth}
    \paragraph{Richiesta (b)}
    Derivarne analiticamente la funzione di trasferimento $A_{BP}(s)=V_{BP}(s)/V_S(s)$ nell'ipotesi
    che l'operazionale sia ideale ed assumendo per semplicit\'a che $R_1=R_2=R=10 kW$, $R3=2R$ e
    $C_1=C_2=C=10 nF$ \textit{
      (suggerimento: si determini la tensione nel punto A del circuito in termini
      di $V_S$ e $V_{BP}$, assumendo che l'operazionale operi in regime lineare come risulta dalle
      precedenti misure; indipendentemente si esprima $V_A$ in termini di $V_{BP}$; si eguaglino
      infine le espressioni ottenute per ottenere la funzione di trasferimento).
    }
  \end{minipage}
}
\paragraph{Risposta (b)}
\begin{equation*}
  \frac{V_S - V_A}{R} = \frac{V_A}{R} + sC(V_A - V_-) + sC(V_A - V_{BP})
\end{equation*}

\begin{equation*}
  sC(V_A - V_-) = \frac{V_- - V_{BP}}{2R}
\end{equation*}
per il principio di corto-circuito virtuale, siccome l'ingresso \underline{non}-invertente viene messo
a terra $V_+ = 0$
\begin{equation*}
  \begin{gathered}
    V_+ = V_- \implies sCV_A = -\frac{V_{BP}}{2R} \\ \implies
    V_A = -\frac{V_{BP}}{2sRC}
  \end{gathered}
\end{equation*}
ricordandoci delle conseguenze del principio di corto-circuito virtuale, l'espressione
nodale iniziale
\begin{equation*}
  V_S - V_A = V_A + sRCV_A + sRC(V_A - V_{BP})
\end{equation*}
sostituendo $V_A$
\begin{equation*}
  V_S + \frac{V_{BP}}{2sRC} = -\frac{V_{BP}}{2sRC} - sRC\frac{V_{BP}}{2sRC} - sRC\left(\frac{V_{BP}}{2sRC} + V_{BP}\right)
\end{equation*}
raccogliendo
\begin{equation*}
  V_S = -\left[\frac{1}{sRC} + 1 + sRC\right]V_{BP}
\end{equation*}
\begin{equation}
  \label{eq:notch-bp-sottrattore-tfunc-bp}
  \colorbox{blu}{\boxed{H_1(s) = -\frac{RCs}{1 + RCs + R^2C^2s^2}}}
\end{equation}

\noindent\colorbox{orange}{
  \begin{minipage}{1.0\linewidth}
    \paragraph{Richiesta (c)}
    Dalla restrizione di $A_{BP}$ sul dominio delle frequenze reali dedurre i valori attesi per i
    parametri del filtro (frequenza centrale, guadagno massimo e larghezza banda) e
    confrontarli con le misure di cui al punto precedente.
  \end{minipage}
}
\begin{equation*}
  \begin{cases}
    \mathcal{R}_{\omega}^N = 0 \\
    \mathcal{I}_{\omega}^N = RC\omega \\
    \mathcal{R}_{\omega}^D = 1 - R^2C^2\omega^2 \\
    \mathcal{I}_{\omega}^D = RC\omega
  \end{cases}
\end{equation*}
il guadagno
\begin{equation*}
  G_{BP}(\omega) = \frac{RC\omega}{\sqrt{(1 - R^2C^2\omega^2)^2 + (RC\omega)^2}}
\end{equation*}
per determinare il guadagno massimo
\begin{equation*}
  G_{BP}(\omega) = \frac{1}{\sqrt{1 + \left(\frac{1 - R^2C^2\omega^2}{RC\omega}\right)^2}}
\end{equation*}
guadagno massimo
\begin{equation}
  \label{eq:notch-bp-sottrattore-gain-max}
  \colorbox{blu}{\boxed{\max_{\omega \in [0, +\infty)} G_{BP}(\omega) = 1}}
\end{equation}
frequenza centrale
\begin{equation}
  \label{eq:notch-bp-sottrattore-frequenza-centrale}
  \colorbox{blu}{\boxed{f = \frac{1}{2\pi RC}}}
\end{equation}
frequenza di taglio
\begin{equation*}
  \begin{gathered}
    \frac{RC\omega}{\sqrt{(1 - R^2C^2\omega^2)^2 + (RC\omega)^2}} = \frac{1}{\sqrt2} \\ \implies
    (RC\omega)^2 = (1 - R^2C^2\omega^2)^2 \implies
    \begin{cases}
      +RC\omega_+ = 1 - R^2C^2\omega_+^2 \\
      -RC\omega_- = 1 - R^2C^2\omega_-^2
    \end{cases} \\ \implies
    \omega_+ = \frac{-1 \pm \sqrt{5}}{2RC} \quad \wedge \quad \omega_- = \frac{1 \pm \sqrt{5}}{2RC}
  \end{gathered}
\end{equation*}
una pulsazione ammissibile deve essere positiva, di conseguenza
\begin{equation*}
  \omega_{L, H} = \frac{\sqrt{5} \pm 1}{2RC}
\end{equation*}
attraverso la sostituzione $\omega \to 2\pi f$ si ottiene la larghezza di banda
\begin{equation}
  \label{eq:notch-bp-sottrattore-bw}
  \colorbox{blu}{\boxed{\Delta f = \frac{1}{2\pi RC}}}
\end{equation}

\noindent\colorbox{orange}{
  \begin{minipage}{1.0\linewidth}
    \paragraph{Richiesta (d)}
    Scrivere la funzione di trasferimento del blocco 2 (ovvero un'espressione di $V_{OUT}$ in
    termini di $V_S$ e $V_{BP}$) e ricavare infine la funzione di trasferimento complessiva del circuito
    ($A_{OUT}(s)=V_{OUT}(s)/V_S(s)$). Che tipo di circuito implementa il blocco 2?
  \end{minipage}
} \smallskip \\
Facendo uso del principio di sovvrapposizione per il quale
\begin{equation*}
  V_{OUT} = H_2(s)\Big|_{V_S = 0}V_{BP} + H_2(s)\Big|_{V_{BP} = 0}V_S
\end{equation*}
\begin{itemize}
\item Per $V_S = 0$
  \begin{equation*}
    \frac{V_{BP} - V_-}{R} = \frac{V_-}{R} + \frac{V_- - V_{OUT}}{R}
  \end{equation*}
  per il principio di corto-circuito virtuale dal fatto che $V_+ = 0$
  \begin{equation*}
    V_+ = V_- \implies \frac{V_{BP}}{R} = \frac{-V_{OUT}}{R}
  \end{equation*}
  da cui
  \begin{equation*}
    \boxed{H_2(s)\Big|_{V_S = 0} = -1}
  \end{equation*}
\item Per $V_{BP} = 0$
  \begin{equation*}
    \frac{-V_-}{R} = \frac{V_- - V_S}{R} + \frac{V_- - V_{OUT}}{R}
  \end{equation*}
  per il principio di corto-circuito virtuale dal fatto che $V_+ = 0$
  \begin{equation*}
    V_+ = V_- \implies 0 = \frac{-V_S}{R} + \frac{-V_{OUT}}{R}
  \end{equation*}
  da cui
  \begin{equation*}
    \boxed{H_2(s)\Big|_{V_{BP} = 0} = -1}
  \end{equation*}
\end{itemize}
infine
\begin{equation}
  \label{eq:notch-bp-sottrattore-tfunc-sum}
  \colorbox{blu}{\boxed{V_{OUT} = V_{BP} + V_S}}
\end{equation}
Il blocco 2 implementa un circuito sommatore. Inoltre siccome
$V_{BP}$ ha dipendenza da $V_S$ attraverso \ref{eq:notch-bp-sottrattore-tfunc-bp}
\begin{equation*}
  \begin{gathered}
    V_{OUT} = H_1(s)V_S + V_S = (1 + H_1(s))V_S \\ \implies
    H(s) \equiv 1 + H_1(s)
  \end{gathered}
\end{equation*}
la funzione di trasferimento complessiva
\begin{equation}
  \label{eq:notch-bp-sottrattore-tfunc}
  \colorbox{blu}{\boxed{H(s) = \frac{1 + R^2C^2s^2}{1 + RCs + R^2C^2s^2}}}
\end{equation}

\noindent\colorbox{orange}{
  \begin{minipage}{1.0\linewidth}
    \paragraph{Richiesta (e)}
    Calcolare i valori attesi del guadagno complessivo a frequenza $f_0$ ed ai limiti del dominio
    di frequenza ($f \to 0$ ed $f \to \infty$) e confrontarli con quanto misurato al punto 1f.
  \end{minipage}
} \smallskip \\
Il guadagno della funzione di trasferimento complessiva \ref{eq:notch-bp-sottrattore-tfunc}
\begin{equation*}
  G(\omega) = \frac{1 - R^2C^2\omega^2}{\sqrt{(1 - R^2C^2\omega^2)^2 + (RC\omega)^2}}
\end{equation*}
i valori attesi
\begin{itemize}
\item
  Per la frequenza $f_0 = 1/(2\pi RC)$ ottenuta da \ref{eq:notch-bp-sottrattore-frequenza-centrale},
  attraverso la sostituzione $\omega_0 \to f_0 / (2\pi)$, si ottiene il guadagno $\colorbox{blu}{\boxed{G(\omega_0) = G(1/RC) = 0}}$.
\item
  Per la frequenza $f \to 0$ il numeratore e denominatore si semplificano $\colorbox{blu}{\boxed{G(0) = 1}}$.
\item
  Per la frequenza $f \to \infty$ il termine dominante del numeratore \'e la potenza quadratica
  e sotto la radice il termine quartico, asintoticamente $\colorbox{blu}{\boxed{G(\omega \to \infty) = 1}}$.
\end{itemize}
Il risultato teorico \'e in accordo con i risultati precedenti, ora devono essere ripetuti i conti per la frequenza
centrale misurata.
